\section{Positive causal experiments}\label{sec:experiments}
\subsection{Types and primitive data}
A standard Borel space is a measurable space isomorphic to a Borel subset of
a Polish space. All primitive spaces below are standard Borel. Time is
initially $t=0,\ldots,N$, with the countable-time extension stated
separately. The unknown parameter is $\lambda\in\Lambda$; it is not a
random variable until a prior is supplied. The hidden state, applied action,
and next report have spaces $X_t,A_t,Y_{t+1}$. A visible history is
\[
 h_t=(a_0,y_1,\ldots,a_{t-1},y_t)\in H_t.
\]
A public initial report, clock increments, and observed cost records can be
included in this notation. An action-availability rule, when present, is
part of the interface. We usually write a common action space to avoid
encoding that rule repeatedly.

\begin{definition}[Raw experiment]\label{def:experiment}
A raw experiment consists of a parameter space $\Lambda$, preparations
$\mu_\lambda\in\cP(X_0)$, and jointly measurable probability kernels
\begin{equation}\label{eq:raw-kernel}
 K_t^\lambda(dx',dy,dc\mid x,a)
\end{equation}
on $X_{t+1}\times Y_{t+1}\times\mathcal C$, together with a class of
executable future tests and a resource interface. The nonnegative vector
$c$ records charged increments, such as raw preparations or physical time.
When it is observed, it is included in $y$; when it is not observed, it is
integrated in the visible law but is still counted in the resource account.
Every kernel has total mass one. A failed measurement and a terminating
measurement are reports, not missing mass.
\end{definition}

The following six requirements will be referred to as the foundational
axioms. They specify the admissible language rather than analytic conclusions.

\begin{description}
\item[A0 (measurable interfaces).] Primitive spaces are standard Borel;
histories and continuation availability are specified.
\item[A1 (positive experiment).] Preparation and transition--observation
kernels are nonnegative and normalized as in \eqref{eq:raw-kernel}.
\item[A2 (causality).] The next action uses only the visible past and
permitted independent randomness. Stopping is nonanticipative.
\item[A3 (executable tests).] The designated future tests have actual
protocols and are closed under the one-step extensions needed for prediction.
\item[A4 (preparation and accounting).] Average risks and acquisition laws
refer to a declared preparation and exploration strategy. Failures, pilots,
and accepted-event probabilities are not removed from their costs.
\item[A5 (charged implementation).] Persistent labels, numerical precision,
workspace, read-only calibration, and simulator state have separate accounts.
\end{description}

In particular these axioms do not assert strict positivity of each report,
finite-dimensionality, stationarity, differentiability, a spectral gap, or a
large-deviation principle. If the physical state is not Markov in the
coordinates initially proposed, its entire necessary past can be included in
$X_t$. This does not prove closure for a smaller observable state.

\subsection{Strategies, paths, and stopping}
A strategy is a sequence of measurable kernels
$\alpha_t(da\mid h_t)$ supported on available actions. It cannot depend on
the unknown true parameter. A finite-state strategy may have an additional
explicit internal state; an unrestricted history strategy need not satisfy a
finite-memory budget.

\begin{proposition}[Path construction]\label{prop:paths}
For every parameter and causal strategy there is a unique consistent
finite-dimensional path law with iterated-kernel expression
\begin{equation}\label{eq:path-law}
 \mu_\lambda(dx_0)\prod_{t=0}^{N-1}
 \alpha_t(da_t\mid h_t)
 K_t^\lambda(dx_{t+1},dy_{t+1},dc_{t+1}\mid x_t,a_t).
\end{equation}
For countably many discrete times these laws extend uniquely to the product
path sigma field. Deterministic evolution and deterministic observations
are included by taking Dirac kernels.
\end{proposition}
\begin{proof}
For finite $N$, integrate bounded measurable cylinder functions from the
last variable backwards. At each step measurability follows from the kernel
property, positivity is preserved, and the integral of one is one. The
resulting probability measures have consistent marginals because the last
kernel integrates to one. The standard kernel-extension theorem gives the
countable product law; this is the Ionescu--Tulcea construction, as treated
in \cite{Kallenberg2021}. Uniqueness follows first on cylinder events and
then on their generated sigma field. A Dirac kernel makes the corresponding
integral substitution along the deterministic trajectory, without changing
the remaining preparation or measurement randomness.
\end{proof}

A bounded stopping time is implemented by an absorbing stopped symbol and
zero subsequent increments. Thus the same construction handles optional
termination. For an unbounded stopping time, an almost-surely finite stopped
path can be considered by truncation; quantitative infinite-horizon error
bounds still require a summable bound or a separate limiting argument.
A random number of collisions is not silently substituted for a
deterministic large-deviation speed.

\subsection{Prepared and parameter-uniform questions}
For statistical decision theory we retain the whole family
$(P_\lambda^\alpha)_\lambda$. A decision rule is one rule, independent of
the true parameter. For Bayesian prediction we instead specify a prior
$\eta(d\lambda)$ and use the joint hidden state $(\lambda,x_t)$ with
preparation $\eta(d\lambda)\mu_\lambda(dx_0)$. Fixed known calibration is
permitted, but it is not a substitute for an unknown parameter.

This distinction also applies to infima over encoders. The infimum for a
fixed preparation need not be achieved by the same encoder for all priors.
Likewise, a family of codebooks indexed by a known experimental calibration
is not an estimator of an unknown calibration. Every quantitative theorem
below states which situation is intended.

\subsection{Instruments and evidence}
After a preparation has been fixed, write the hidden state simply as $x$.
For a report event $B\subset Y_{t+1}$, define the positive instrument
\[
 (\mathsf T_{t,a,B}f)(x)
 =\int \1_B(y)f(x')K_t(dx',dy\mid x,a).
\]
It is countably additive in $B$ for nonnegative $f$, positive in $f$, and
$\mathsf T_{t,a,Y}1=1$. With finite reports write
$\mathsf T_{t,a,y}=\mathsf T_{t,a,\{y\}}$.
Starting with $b_0=\mu$, define the unnormalized history functional by
\begin{equation}\label{eq:unnormalized}
 b_{t+1}(f)=b_t(\mathsf T_{t,a_t,y_{t+1}}f),
 \qquad Z_t=b_t(1).
\end{equation}
These functionals do not include the probability of the chosen command.
They are likelihood/evidence functionals conditional on the realized
command sequence; the actual joint history law also includes the strategy
kernels in \eqref{eq:path-law}.

\begin{proposition}[Normalized update]\label{prop:bayes}
For finite reports and $Z_t>0$, let $\pi_t=b_t/Z_t$. If
$q_t(y\mid\pi_t,a)=\pi_t(\mathsf T_{t,a,y}1)>0$, then
\begin{equation}\label{eq:bayes}
 \pi_{t+1}(f)=
 \frac{\pi_t(\mathsf T_{t,a,y}f)}
      {\pi_t(\mathsf T_{t,a,y}1)},
 \qquad
 Z_{t+1}=Z_tq_t(y\mid\pi_t,a).
\end{equation}
The formula remains valid for adaptive actions. When the evidence is zero,
the unnormalized functional is zero and no normalized posterior is specified
by this formula.
\end{proposition}
\begin{proof}
Divide \eqref{eq:unnormalized} by its value at one. Positivity shows that a
positive functional of total mass zero vanishes on every bounded test.
For adaptive actions, condition on the visible history: the strategy's
probability of the current action is fixed by that history and is independent
of the hidden state conditional on it. This factor cancels in the posterior.
An induction in \eqref{eq:path-law} verifies the conditional expectation
identity for every bounded hidden-state test.
\end{proof}

For a non-atomic report space, conditioning is by disintegration of the
joint next-state/report law, not by dividing the mass of a singleton. Under
a fixed preparation a standard Borel conditional distribution exists almost
everywhere. A single Borel update jointly in all proposed states, inputs,
and reports is additional implementation data, or must be justified by an
applicable parametrized disintegration theorem. We will not use a
simultaneous version across a nondominated parameter family without that
justification. All explicit quantitative models below have finite reports
and Borel formulas.

\subsection{Preparing a conditional experiment}
If an event has acceptance probability $p>0$, independently repeating a
preparation until acceptance implements the conditional law. Its preparation
count $T$ satisfies $\Pp(T>n)=(1-p)^n$ and $\E T=1/p$. A cap $B$ instead
creates an abort symbol with probability $(1-p)^B$. This is a different
experiment from exact conditioning. The cap is necessary for a uniform
pathwise preparation budget. No lower bound or implementation claim below
replaces an actual report word by a probability-one conditioning event.
